% Pregunta 8. Sistema bancario, estabilidad financiera y macroprudencial (bloque V.2) · peso bajo
\section*{Pregunta 8. Sistema bancario y estabilidad financiera}

\begin{enunciado}
\textbf{8.} Valore la situación del \textbf{sistema bancario} español en términos de solvencia, calidad de los activos y rentabilidad. ¿Qué riesgos para la \textbf{estabilidad financiera} aprecia? ¿Qué medidas de política macroprudencial considera oportunas?

\smallskip
\textbf{Peso bajo} · \tiempo{13} · guía: V.2 (endeudamiento, indicadores bancarios, riesgos, macroprudencial; semáforo de la banca) · cuadros 4 y 8; gráfico 4.
\end{enunciado}

\begin{examen}
\textit{El sistema financiero español parece resiliente: la banca es muy rentable, tiene poca morosidad y cumple con holgura los requisitos de capital, aunque con menos capital de máxima calidad que la media europea, y la deuda privada está en mínimos. El riesgo que hay que vigilar es la vivienda.}

\rotulo{a) Solvencia, calidad de los activos y rentabilidad}
La \textbf{solvencia} es suficiente, pero no holgada: el capital ordinario de nivel 1 (CET1) es el 13,9~\% de los activos ponderados por riesgo y el capital total, el 18,1~\% (cuadro~8), por encima de los requisitos medios del Mecanismo Único de Supervisión (en torno al 11~\% de CET1), pero por debajo de la media de la UE (16,4~\%). La \textbf{calidad de los activos} es buena: los préstamos dudosos bajan del 8,0~\% en 2014 al 2,4~\%, por debajo del 3~\% que se considera bajo, aunque algo por encima de la media europea (2,0~\%). La \textbf{rentabilidad} es alta: un ROE del 14,2~\% (9,3~\% en la UE), por encima del coste del capital, lejos del $-25{,}6~\%$ de 2012. No se debe a un mayor apalancamiento, sino a un activo más rentable (ROA del 0,96~\% frente al 0,71~\%): el margen de intereses se beneficia de los tipos más altos y de un crédito a tipo variable que se reprecia deprisa.

\rotulo{b) Riesgos para la estabilidad financiera}
El principal es la \textbf{vivienda}: su precio sube un 12,7~\% en 2025, por encima del umbral del 9~\% del procedimiento de desequilibrios, y crece más que la inversión residencial (más precio que cantidad). Lo mitiga que no hay un auge del crédito: la deuda de los hogares está en el 42,9~\% del PIB y el crédito a hogares crece un 4,5~\%. El segundo es el \textbf{endurecimiento monetario} unido a una desaceleración: con mucho crédito a tipo variable, la subida de los tipos se traslada enseguida a las cuotas y podría elevar la morosidad, que va retrasada respecto al ciclo. Y, en las empresas, el encarecimiento de la energía y los aranceles. El bucle bancario-soberano preocupa poco con un diferencial de 45~pb.

\rotulo{c) Medidas macroprudenciales}
En una fase alcista con la vivienda al alza, lo prudente es \textbf{acumular capital ahora}: mantener activado el colchón de capital anticíclico y elevarlo si el crédito se acelera, y moderar el reparto de dividendos para acercar el CET1 a la media europea. Para la vivienda, tener preparados \textbf{límites a los prestatarios} (préstamo sobre el valor de la vivienda y sobre la renta) y activarlos si el crédito hipotecario se dispara. El problema de fondo de la vivienda, la oferta, no lo resuelve la política macroprudencial (pregunta~10).
\end{examen}

\begin{complemento}{8.1 · Por qué se contesta así: el semáforo y el orden de la guía}
\que{Qué se aplica.} Guía, V.2: empezar con una frase que resuma el estado del sistema y seguir el orden \emph{endeudamiento, banca, riesgos y políticas}; las ratios con sus referencias (gráfico~\ref{fig:semaforo}); $\mathrm{ROE}=\mathrm{ROA}\times\text{activo}/\text{patrimonio}$; la vivienda como «principal activo, principal deuda y garantía del crédito»; instrumentos macroprudenciales en la fase alcista.

\begin{figuracaja}
\caption{Semáforo de la banca: referencias de las ratios principales (infografía de la guía)}\label{fig:semaforo}
\grafico[0.58\linewidth]{esquema-semaforo-banca}
\fuente{Guía del dictamen, V.2. España en 2025 (cuadro~8): CET1 13,9~\% (\textcolor{tceedorado}{\textbf{vigilar}}, cerca de holgado); dudosos 2,4~\% (\textcolor{tceeverde}{\textbf{holgado}}); ROE 14,2~\% (\textcolor{tceeverde}{\textbf{holgado}}). La liquidez y la eficiencia no vienen en el cuadernillo.}
\end{figuracaja}

\que{Cálculos} (cuadro~8): apalancamiento implícito, $\mathrm{ROE}/\mathrm{ROA}=14{,}2/0{,}96\approx14{,}8$ (el patrimonio es el 6,8~\% del activo); en la UE, $9{,}3/0{,}71\approx13{,}2$. España apalanca algo más, pero la diferencia de ROE (4,9 puntos) se explica sobre todo por el ROA (0,96 frente a 0,71). Margen de CET1 sobre los requisitos medios del MUS (11,2~\%, guía V.2 \parencite{BCE2025SREP}): unos 2,7 puntos.

\que{Más información de la que trae el cuadernillo.} El colchón de capital anticíclico de España es del 1~\% desde octubre de 2026 \parencite{BdE2025CCA} (guía, tabla de Basilea~III), y las estadísticas supervisoras del Banco de España daban a mediados de 2026 un ROE del 16~\% y un CET1 del 14~\% \parencite{BdE2026EstadSupervisoras}. Con ese dato, la recomendación de (c) es \emph{mantenerlo y estar listos para subirlo}, no activarlo.

\que{Errores típicos.} Confundir el ROA bancario con el de una empresa no financiera; decir que la banca «no cumple» por tener menos CET1 que la media (cumple con margen); recomendar medidas de crisis (liberar colchones) en una fase alcista. \que{Qué añadir con más tiempo.} La brecha crédito-PIB (no viene; el flujo de crédito del 2,6~\% del PIB sugiere que es negativa, guía V.2) y la liquidez (LCR, créditos sobre depósitos). \que{Qué recortar en el examen.} Nada: es la otra pregunta corta.
\end{complemento}
